%%%%%%%%%%%%%%%%%%%%%%%%%%%%%%%%%%%%%%%%%%%%%%%%%%%%%%%%%%%%%%%%%%%
%%%                                                             %%%
%%%      File: poster1.tex                                      %%%
%%%    Author: Sabine Schulte im Walde                          %%%
%%%   Purpose: Research Seminar                                 %%%
%%%                                                             %%%
%%%%%%%%%%%%%%%%%%%%%%%%%%%%%%%%%%%%%%%%%%%%%%%%%%%%%%%%%%%%%%%%%%%


% a0poster Portrait Poster
% LaTeX Template
% Version 1.0 (22/06/13)
%
% The a0poster class was created by:
% Gerlinde Kettl and Matthias Weiser (tex@kettl.de)
% 
% This template has been downloaded from:
% http://www.LaTeXTemplates.com
%
% License:
% CC BY-NC-SA 3.0 (http://creativecommons.org/licenses/by-nc-sa/3.0/)


%----------------------------------------------------------------------------------------
%	PACKAGES AND OTHER DOCUMENT CONFIGURATIONS
%----------------------------------------------------------------------------------------

\documentclass[a0,portrait]{a0poster}

\usepackage{multicol} % This is so we can have multiple columns of text side-by-side
\columnsep=100pt      % This is the amount of white space between the columns in the poster
\columnseprule=3pt    % This is the thickness of the black line between the columns in the poster

\usepackage[svgnames]{xcolor} % Specify colors by their
                              % 'svgnames'. For a full list of all
                              % colors available see here:
                              % http://www.latextemplates.com/svgnames-colors

\usepackage{times} % Use the times font

\usepackage{graphicx}     % Required for including images
\graphicspath{{figures/}} % Location of the graphics files
\usepackage{booktabs}     % Top and bottom rules for table
\usepackage[font=small,labelfont=bf]{caption}   % Required for specifying captions to tables and figures
\usepackage{amsfonts, amsmath, amsthm, amssymb} % Required for math fonts, symbols and environments
\usepackage{wrapfig}      % Allows wrapping text around tables and figures


\begin{document}

%----------------------------------------------------------------------------------------
%	POSTER HEADER 
%----------------------------------------------------------------------------------------

% The header is divided into two boxes:
% The first is 75% wide and houses the title, subtitle, names, university/organization and contact information
% The second is 25% wide and houses a logo for your university/organization or a photo of you
% The widths of these boxes can be easily edited to accommodate your content as you see fit

\begin{minipage}{0.75\linewidth}
  \veryHuge \color{NavyBlue} \textbf{Research Seminar at the IMS} \color{Black} % Title
  \vspace{+1cm}\\
  \Huge \textit{Presentation on paper : GQA: A New Dataset for Real-World Visual Reasoning and Compositional Question Answering. Hudson et al. (2019)}\\[2cm] % Subtitle
  \huge \textbf{Shawon Ashraf}\\[0.5cm] % Author(s)
  \huge IMS, University of Stuttgart\\[0.4cm] % University/Organization
  \Large \texttt{shawon.ashraf@ims.uni-stuttgart.de}\\
\end{minipage}
%
\begin{minipage}{0.25\linewidth}
  \includegraphics[width=15cm]{logo_red_text_left}\\
\end{minipage}

\vspace{5cm} % A bit of extra whitespace between the header and the poster content

%----------------------------------------------------------------------------------------

\begin{multicols}{2} % This is how many columns your poster will be broken into. A portrait poster is generally split into 2 columns.


%----------------------------------------------------------------------------------------
%	INTRODUCTION/MOTIVATION
%----------------------------------------------------------------------------------------

\color{SaddleBrown} % SaddleBrown color for the introduction

\section*{Introduction/Motivation}

Introduce/motivate your poster.
\vspace{+5mm}\\
Do you need citations \cite{Smith:2012qr}?


%----------------------------------------------------------------------------------------
%	OBJECTIVES
%----------------------------------------------------------------------------------------

\color{DarkSlateGray} % DarkSlateGray color for the rest of the content

\vspace{+3cm}
\section*{Main Objectives}

\begin{enumerate}
\item Lorem ipsum dolor sit amet, consectetur.
\item Nullam at mi nisl. Vestibulum est purus, ultricies cursus volutpat sit amet, vestibulum eu.
\item Praesent tortor libero, vulputate quis elementum a, iaculis.
\item Phasellus a quam mauris, non varius mauris. Fusce tristique, enim tempor varius porta, elit purus commodo velit, pretium mattis ligula nisl nec ante.
\item Ut adipiscing accumsan sapien, sit amet pretium.
\item Estibulum est purus, ultricies cursus volutpat
\item Nullam at mi nisl. Vestibulum est purus, ultricies cursus volutpat sit amet, vestibulum eu.
\item Praesent tortor libero, vulputate quis elementum a, iaculis.
\end{enumerate}


%----------------------------------------------------------------------------------------
%	MATERIALS AND METHODS
%----------------------------------------------------------------------------------------

\vspace{+2cm}
\section*{Materials and Methods}

\begin{enumerate}
\item Lorem ipsum dolor sit amet, consectetur.
\item Nullam at mi nisl. Vestibulum est purus, ultricies cursus volutpat sit amet, vestibulum eu.
\item Praesent tortor libero, vulputate quis elementum a, iaculis.
\item Phasellus a quam mauris, non varius mauris. Fusce tristique, enim tempor varius porta, elit purus commodo velit, pretium mattis ligula nisl nec ante.
\item Ut adipiscing accumsan sapien, sit amet pretium.
\item Estibulum est purus, ultricies cursus volutpat
\item Nullam at mi nisl. Vestibulum est purus, ultricies cursus volutpat sit amet, vestibulum eu.
\item Praesent tortor libero, vulputate quis elementum a, iaculis.
\end{enumerate}


%------------------------------------------------

\vspace{+1cm}
\subsection*{Mathematical Section 1}

Nulla vel nisl sed mauris auctor mollis non sed. 

\begin{equation}
E = mc^{2}
\label{eqn:Einstein}
\end{equation}
%
Curabitur mi sem, pulvinar quis aliquam rutrum. (1) edf (2), $\Omega=[-1,1]^3$, maecenas leo est, ornare at. $z=-1$ edf $z=1$ sed interdum felis dapibus sem. $x$ set $y$ ytruem. 
Turpis $j$ amet accumsan enim $y$-lacina; 
ref $k$-viverra nec porttitor $x$-lacina. 

\vspace{+1cm}
\subsection*{Mathematical Section 2}

Vestibulum ac diam a odio tempus congue. Vivamus id enim nisi:

\begin{eqnarray}
\cos\bar{\phi}_k Q_{j,k+1,t} + Q_{j,k+1,x}+\frac{\sin^2\bar{\phi}_k}{T\cos\bar{\phi}_k} Q_{j,k+1} &=&\nonumber\\ 
-\cos\phi_k Q_{j,k,t} + Q_{j,k,x}-\frac{\sin^2\phi_k}{T\cos\phi_k} Q_{j,k}\label{edgek}
\end{eqnarray}
and
\begin{eqnarray}
\cos\bar{\phi}_j Q_{j+1,k,t} + Q_{j+1,k,y}+\frac{\sin^2\bar{\phi}_j}{T\cos\bar{\phi}_j} Q_{j+1,k}&=&\nonumber \\
-\cos\phi_j Q_{j,k,t} + Q_{j,k,y}-\frac{\sin^2\phi_j}{T\cos\phi_j} Q_{j,k}.\label{edgej}
\end{eqnarray} 
%
Nulla sed arcu arcu. Duis et ante gravida orci venenatis tincidunt. Fusce vitae lacinia metus. Pellentesque habitant morbi. $\mathbf{A}\underline{\xi}=\underline{\beta}$ Vim $\underline{\xi}$ enum nidi $3(P+2)^{2}$ lacina. Id feugain $\mathbf{A}$ nun quis; magno.


%----------------------------------------------------------------------------------------
%	RESULTS 
%----------------------------------------------------------------------------------------

\vspace{+2cm}
\section*{Results}

Describe your results.\\
%
\begin{wraptable}{l}{12cm} % Left or right alignment is specified in the first bracket, the width of the table is in the second
  \begin{tabular}{l l l}
    \toprule
    \textbf{Treatments} & \textbf{Response 1} & \textbf{Response 2}\\
    \midrule
    Treatment 1 & 0.0003262 & 0.562 \\
    Treatment 2 & 0.0015681 & 0.910 \\
    Treatment 3 & 0.0009271 & 0.296 \\
    \bottomrule
  \end{tabular}
  \captionof{table}{\color{Green} Table caption}
\end{wraptable}
%
Phasellus imperdiet, tortor vitae congue bibendum, felis enim sagittis lorem, et volutpat ante orci sagittis mi. Morbi rutrum laoreet semper. Morbi accumsan enim nec tortor consectetur non commodo nisi sollicitudin. Proin sollicitudin. Pellentesque eget orci eros. Fusce ultricies, tellus et pellentesque fringilla, ante massa luctus libero, quis tristique purus urna nec nibh. Fusce tristique, enim tempor varius porta, elit purus commodo velit, pretium mattis ligula nisl nec ante.

\columnbreak

In hac habitasse platea dictumst. Etiam placerat, risus ac.

\begin{center}\vspace{1cm}
  \includegraphics[width=0.6\linewidth]{placeholder}
  \captionof{figure}{\color{Green} Figure caption}
\end{center}\vspace{1cm}

Adipiscing lectus in magna blandit:

\begin{center}\vspace{1cm}
  \begin{tabular}{l l l l}
    \toprule
    \textbf{Treatments} & \textbf{Response 1} & \textbf{Response 2} \\
    \midrule
    Treatment 1 & 0.0003262 & 0.562 \\
    Treatment 2 & 0.0015681 & 0.910 \\
    Treatment 3 & 0.0009271 & 0.296 \\
    \bottomrule
  \end{tabular}
  \captionof{table}{\color{Green} Table caption}
\end{center}\vspace{1cm}
%
Vivamus sed nibh ac metus tristique tristique a vitae ante. Sed lobortis mi ut arcu fringilla et adipiscing ligula rutrum. Aenean turpis velit, placerat eget tincidunt nec, ornare in nisl. In placerat.

\begin{center}\vspace{1cm}
  \includegraphics[width=0.6\linewidth]{placeholder}
  \captionof{figure}{\color{Green} Figure caption}
\end{center}\vspace{1cm}


%----------------------------------------------------------------------------------------
%	CONCLUSIONS
%----------------------------------------------------------------------------------------

\color{SaddleBrown} % SaddleBrown color for the conclusions to make them stand out

\vspace{+2cm}
\section*{Conclusions}

\begin{itemize}
\item Pellentesque eget orci eros. Fusce ultricies, tellus et pellentesque fringilla, ante massa luctus libero, quis tristique purus urna nec nibh. Phasellus fermentum rutrum elementum. Nam quis justo lectus.
\item Vestibulum sem ante, hendrerit a gravida ac, blandit quis magna.
\item Donec sem metus, facilisis at condimentum eget, vehicula ut massa. Morbi consequat, diam sed convallis tincidunt, arcu nunc.
\item Nunc at convallis urna. isus ante. Pellentesque condimentum dui. Etiam sagittis purus non tellus tempor volutpat. Donec et dui non massa tristique adipiscing.
\end{itemize}
<
\color{DarkSlateGray} % Set the color back to DarkSlateGray for the rest of the content


%----------------------------------------------------------------------------------------
%	FORTHCOMING RESEARCH
%----------------------------------------------------------------------------------------

\vspace{+3cm}
\section*{Forthcoming Research}

Vivamus molestie, risus tempor vehicula mattis, libero arcu volutpat purus, sed blandit sem nibh eget turpis. Maecenas rutrum dui blandit lorem vulputate gravida. Praesent venenatis mi vel lorem tempor at varius diam sagittis. Nam eu leo id turpis interdum luctus a sed augue. Nam tellus.


%----------------------------------------------------------------------------------------
%	REFERENCES
%----------------------------------------------------------------------------------------

\vspace{+3cm}
\nocite{*} % Print all references regardless of whether they were cited in the poster or not
\bibliographystyle{plain} % Plain referencing style
\bibliography{sample1} % Use the example bibliography file sample.bib


%----------------------------------------------------------------------------------------
%	ACKNOWLEDGEMENTS
%----------------------------------------------------------------------------------------

\vspace{+3cm}
\section*{Acknowledgements}

Etiam fermentum, arcu ut gravida fringilla, dolor arcu laoreet justo, ut imperdiet urna arcu a arcu. Donec nec ante a dui tempus consectetur. Cras nisi turpis, dapibus sit amet mattis sed, laoreet.


%----------------------------------------------------------------------------------------

\end{multicols}

\end{document}


%% --- END OF FILE
